\section{Conclusion}

The context window of a large language model is not memory. It is
L1~cache---a small, fast, expensive level in what should be a
managed hierarchy. The field treats it as the entire memory system.
There is no L2, no paging, no eviction policy, no working set
estimation. The result: 21.8\% structural waste in 4.45~billion
measured tokens, with the model reprocessing stale content at
84.4$\times$ amplification.

\textsc{Pichay} demonstrates that the virtual memory abstractions
developed for physical memory in the 1960s apply directly to LLM
context windows. Demand paging with a simple FIFO policy recovers
36~percentage points of context capacity in steady-state use.
Fault-driven pinning reduces repeated faults on working-set content in
steady-state use.
Retrieval handles---designed as space-saving markers---function as
late-binding anchors that models understand without instruction.
Graduated pressure zones and cooperative cleanup tags give the model
agency over its own memory management---a capability with no analogue
in hardware memory hierarchies.

The system requires no changes to models, clients, or inference
APIs. It is implemented as a transparent proxy and deployed in
production use. The paper you are reading was written through it.

The deeper contribution is the architectural observation: LLM
systems need not one large context window but a hierarchy of
managed levels---from the generation window (L1) through persistent
cross-session memory (L4)---connected by eviction policies and
fault mechanisms. The first three levels are deployed; the first
two are evaluated. L3---model-initiated conversation compaction---is
the newest mechanism, and its evaluation is ongoing. The solutions
for the remaining level exist in the OS literature.
What remains is to connect them.

% ====================================================================
